\section*{The Propers: Detailed Commentary}

\subsection*{Four opening voices: confidence, petition, lament, acclamation\enspace\properrefs{Int., Coll., Grad., All.}}

These four texts do not repeat one mood. The Introit makes confession through assertions and questions; the Collect asks through coordinated clauses; the Gradual rearranges a communal lament into plea and versicle; the Alleluia acclaims judgment and refuge before the Gospel. Their formal differences keep the Sunday's dependence upon God from becoming a generic theme.

\begin{fourstable}{Proper}{Textual form}{Complete-context depth}{Inherited reading}
Introit & The Lord \emph{is} light, salvation, and defender; ``whom?'' and ``from whom?'' deny the enemy the final power to name reality. & Psalm 26 moves from threatened confidence toward one desire: God's dwelling, presence, and patient waiting. The appointed army verse intensifies danger rather than pretending it absent. & Augustine, \emph{Psalm 26}, 1--3, directly reads received illumination and salvation as the ground of courage under assault. He does not turn confidence into immunity from persecution.\\
Collect & One petition coordinates two subjects and effects: the world's moving \latin{cursus} is to be directed peacefully by God's order, and God's Church is to rejoice in tranquil devotion. & The prayer asks neither that history stop nor that ecclesial joy be detached from the world's continuing motion. Its peace is received direction, not mere stillness. & Augustine, \emph{City of God} XIX.13, names peace the tranquility of order and distinguishes temporal peace from final peace; the checked witnesses supply no secure direct commentary on this Collect.\\
Gradual & The proper first joins \latin{Propitius esto} to the nations' question, then returns in the versicle to \latin{Adiuva nos} and \latin{libera nos}. Pardon, God's name, help, and deliverance remain one communal plea. & Psalm 78 speaks after Jerusalem and the Temple have been devastated. The hostile question ``Where is their God?'' threatens the nations with despair as well as Israel with shame. & Augustine, \emph{Psalm 78}, 12--13, directly calls God the healer, holds grace and the helped person's action together, and reads concern for the nations within the petition.\\
Alleluia & Enthroned equity and refuge for the poor are acclaimed in one verse rather than balanced as rival divine attributes. & Psalm 9 moves through judgment of oppressors toward the poor person's refuge and hope in God's name; refuge belongs inside just judgment, not outside it. & Augustine, \emph{Psalm 9}, 6 and 9--10, directly distinguishes God's knowledge of hearts from human judgment and locates refuge in him. Neither the psalm nor Augustine identifies the throne with Christ's boat-seat.\\
\end{fourstable}

The sequence consequently deepens rather than cancels its opening confidence. Those who name God as light still ask that the world be directed; those who ask pardon still acclaim the Judge; those who acclaim judgment approach that Judge as refuge.

\sectionguard
\subsection*{Romans 8: travail tends toward bodily glory\enspace\properrefs{Ep.}}

Romans 8:18--23 is not a general maxim that pain eventually pays. Paul compares present sufferings with glory to be revealed, then gives creation a voice: it was subjected to vanity in hope, groans in travail, and awaits the revelation of God's children. Believers possess the firstfruits of the Spirit while still groaning for adoption's completion, ``the redemption of our body.'' Possession and waiting, gift and ache, remain simultaneous.

\begin{fourstable}{Pauline hinge}{Chrysostom's direct exegesis}{Doctrinal force}{Precise sense}
Expectation under hope & Paul personifies creation, subjected to corruption in connection with fallen humanity yet promised liberation. & Non-rational creation is good and belongs to humanity's history and future honor. & Personification does not assign human psychology to every creature; creation and humanity remain distinct.\\
Groaning and travail & Chrysostom hears labor pains ordered toward newness, not pain circling without an end. & Christian lament can be unguardedly truthful because hope rests in God's promised act. & Birth imagery does not make suffering good in itself.\\
Firstfruits and bodily redemption & The Spirit is a genuine pledge while believers await the body's complete, irreversible freedom from corruption. & Already and not-yet stay joined; salvation is not escape from embodiment. & Firstfruits rule out despair, waiting rules out presumption, and resurrection exceeds return to present mortality.\\
\end{fourstable}

\begin{sourcecard}{St. Irenaeus --- the mortal body vivified}
\emph{Against Heresies} V.7.1--2 appeals to Romans 8:11: the mortal body that dies is the body the Spirit vivifies. ``Spiritual body'' therefore names this body raised through the Spirit for enduring life, not a replacement for discarded matter. Irenaeus comments on a closely related Pauline verse, not directly on Romans 8:18--23.
\end{sourcecard}

\sectionguard
\subsection*{Luke 5: Christ teaches, commands, gives, and calls\enspace\properrefs{Gosp.}}

Luke's narrative order prevents the catch from becoming a free-standing success story. The crowd first presses to hear God's word; Jesus teaches from Simon's boat; only then does he command another cast. Simon reports expert failure without disguise and obeys before the result. The catch exceeds private capacity, Peter confesses sin, Christ answers fear, and the fishermen leave the abundance to follow its giver.

\begin{fourstable}{Narrative turn}{St. Cyril of Alexandria}{St. Ambrose}{Consequence}
Word before sign & The future fishers are first caught by Christ's teaching; the sign pledges the teacher's almighty power. & Christ chooses Peter's boat as a place of teaching before directing its labor. & Mission receives doctrine and authority; results do not create them.\\
\latin{Duc in altum} & Creation obeys Christ's will, and he supplies what the net cannot produce. & In IV.71 the ``deep'' includes the depth of confession concerning the Son's divine generation. & Launching deep is not merely motivational daring.\\
The catch & Cyril reads it as a type of apostolic preaching among the nations; Christ continues to fill the net. & Apostolic discourse brings persons from depths and obscurity into spiritual light alive. & The type does not promise numerical abundance to every faithful undertaking.\\
Partners and boats & Later pastors cooperate in a work whose efficacy remains Christ's. & Ambrose permits several ecclesial readings of the two boats or peoples. & Cooperation is real; no one symbolic identification exhausts the scene.\\
Confession and call & Peter's fear is reverence before Christ's purity and power. & Empty presumption yields to humility, confession, gift, and vocation. & Grace answers guilt without pretending Peter is already complete.\\
\end{fourstable}

\sectionguard
\subsection*{Luke's burdened catch and John's completed shore\enspace\properrefs{Gosp.}}

Augustine reads Luke 5 beside John 21 without denying either event's literal history. In \emph{Tractate 122 on John}, 7, the differences form an ecclesiological contrast: Luke's catch signifies the Church in the present age, while John's post-Resurrection catch signifies her final condition.

\begin{fourstable}{Feature}{Luke 5}{John 21}{Augustine's reading}
Setting and catch & Before the Passion; immense and unnumbered; boats remain on the water. & After the Resurrection; 153 fish are drawn to Christ on shore. & Present gathering precedes resurrection's stability and the number completed by God.\\
Net & It strains and tears. & It is explicitly not torn. & Present communion suffers rupture; final communion cannot be broken.\\
Fish & No right-side cast or final sorting is narrated. & The cast is on the right and the fish reach shore. & The present gathering is mixed; final manifestation belongs to the end.\\
Boats & Two boats fill and nearly sink. & The catch reaches land without sinking the boat. & Ecclesial burden does not mean Christ's word has failed.\\
\end{fourstable}

The contrast guards both directions: abundance cannot certify every person or structure, and strain cannot refute Christ's agency. Purity and completion belong to the shore still awaited.

St.~Anthony of Padua's medieval sermon for the Fifth Sunday after Pentecost preaches on the same catch. He repeatedly returns to Peter's \latin{in verbo tuo}: nets cast on the preacher's own word catch nothing or only the praise-seeking croak of a frog, while fruit received on Christ's word must be attributed to grace. At Peter's knees he holds fear over sin together with hope in the Redeemer's mercy, then treats leaving all as freedom for following. Cyril emphasizes the Lord's continuing gift, Ambrose the apostolic and ecclesial depths of the boat, and Anthony the preacher's temptation to appropriate the catch; together they give abundance an ascetical as well as ecclesial test.

\sectionguard
\subsection*{The offered petition is for wakeful interior sight\enspace\properrefs{Off.}}

Psalm 12 moves from delay and sorrow through a plea for illuminated eyes to trust in mercy. The Offertory appoints the petition, not a boast about achieved clarity.

\begin{sourcecard}{St. Augustine --- on Psalm 12:4--5}
In sections 3--6 Augustine identifies the enemy as the devil or entrenched carnal habit. The eyes are the eyes of the heart, liable to a pleasurable eclipse through sin; sleep in death names the danger of losing faith's light. Even steadfastness is attributed to God's mercy, so the petition opposes both despair and self-reliance.
\end{sourcecard}

At Offertory the gift-bearer still asks to see, remain awake, and escape the enemy's claim. The Introit's indicative---the Lord is my light---and this imperative---illuminate my eyes---remain distinct forms of dependence.

\sectionguard
\subsection*{The Secret asks acceptance and merciful compulsion\enspace\properrefs{Sec.}}

Three clauses carry the Secret's argument; the hand missal prints the complete prayer.

\begin{threestable}{Latin hinge}{Grammatical force}{Theological pressure}
\latin{Oblationibus nostris ... susceptis} & Ablative absolute: our oblations having been received & Reception is God's act; the gifts are truly ours without becoming leverage over him.\\
\latin{placare ... et ad te} & A passive imperative asks favorable disposition; \latin{et} coordinates a second petition whose end is God. & Sacrifice is ordered toward reconciliation, while conversion is personal direction toward God.\\
\latin{nostras etiam rebelles compelle propitius voluntates} & The wills are the object; resistance is named; efficacious action is requested in a favorable, merciful mode. & The prayer neither conceals resistance nor grants it the final word; divine action is not human coercion or annihilation of agency.\\
\end{threestable}

\begin{studybox}{Two coordinated petitions}
The prayer coordinates two petitions: accept the oblations and be favorable; mercifully compel rebellious wills toward God. It does \emph{not} identify the wills as the oblations. Their relation within one prayer is liturgically suggestive, but the Latin does not say ``the will as interior offering.''
\end{studybox}

Augustine's \emph{Confessions} VIII.5 and 8--12 narrates a will divided through entrenched habit without dividing the person into alien substances. \emph{Grace and Free Will} 31--33 says that God gives what he commands by helping the commanded person obey; grace initiates willing and cooperates in action. Chrysostom, on Philippians 2:12--13, likewise joins God's working of willing and action to exhortation, effort, and obedience. These sources illuminate the doctrine at stake; none comments on this Secret.

\sectionguard
\subsection*{Refuge received, purification still requested\enspace\properrefs{Comm., Postcomm.}}

Psalm 17 is David's song of deliverance, and Augustine reads its speaking subject as whole Christ, head and body. In sections 2--3, ``my strength'' means the Lord by whom the speaker is strong; refuge follows divine deliverance; hope rests on the help of God's prior call; the defender prevents proud reliance upon self. Communion therefore gives the Church a voice of received strength rather than acquired invulnerability.

The Postcommunion preserves the same dependence after reception. The mysteries are \latin{sumpta}, taken or received, and precisely as received they are asked to purify and protect. The prayer names no visible result as proof that their work is finished.

\begin{sourcecard}{St. Irenaeus --- Eucharist and the body's future}
\emph{Against Heresies} V.2.2--3 argues that Eucharistic bread and cup nourish the flesh with Christ's body and blood; that same flesh is capable of God's gift and will rise incorruptibly. Irenaeus's doctrine closes the distance between Romans' bodily redemption and sacramental reception, though he does not comment on this Communion antiphon or Postcommunion.
\end{sourcecard}

\sectionguard
\section*{The Propers: Notable and Quotable}

{\small
\begin{itemize}[itemsep=0.35em,topsep=0.2em,parsep=0pt]
\item \textbf{\latin{Dominus illuminatio mea} --- University of Oxford motto and a Vice-Chancellor's pun.} The Introit's ``Lord is my light'' has served on Oxford's arms since at least the later sixteenth century, shifting psalmic trust into institutional claims about learning. In the 2023 Oration, Irene Tracey deliberately dropped letters to make \latin{Domus illuminat nos}---``our home enlightens us''---and redirected divine illumination toward collegiate community. The joke is her update, not a translation of the psalm.

\item \textbf{\latin{Duc in altum} --- Mitt Romney's Southern Virginia University commencement, 2013.} After quoting Luke, Romney turned fishing depth into a whole-life metaphor and urged graduates to marry, start families, find God, pursue careers, and serve. A concrete maneuver after an empty night's work becomes unusually specific commencement advice; the promise of filled nets belongs to the speech's rhetoric, not to Luke as a guarantee of worldly success.

\item \textbf{``The whole creation groaneth'' --- William Lyon Phelps, \emph{Reading the Bible} (1919), pp.~8--9.} Phelps remembered picturing everyone in the world suffering green-apple indigestion at the same instant and drawing them in a circle. Yale president Arthur Hadley later supplied the punchline: unauthorized apple-eating really did make creation groan. Paul's cosmic travail becomes gastrointestinal slapstick and an Eden pun, not an alternative exegesis.
\end{itemize}
}

\sectionguard
\section*{The Propers: Interpretive Possibilities}

\noindent\emph{The following exploratory proposals were developed in an AI-assisted editorial process. They are attributed to none of the cited authorities and claim no historical compositional intent. Each states its anchors and limits.}

\Needspace{8\baselineskip}
\noindent\textbf{Peace as a received direction.\enspace\properrefs{Coll., Gosp., Postcomm.}}
The Collect asks not that the \latin{mundi cursus} cease, but that it be peacefully directed; the Gospel answers with a chain of received directions---put out, lower, gather, follow---and the Postcommunion ends by asking that the people be protected. The connecting mechanism is directed movement whose source and safeguard lie beyond the mover. Read element by element, peace can sound static, obedience episodic, and protection merely defensive; together they suggest that Christian peace may be the stable orientation of a life still under way. The pastoral fruit is a sharper test for discernment: neither motion nor repose is decisive until its giver, direction, and end are named. The controlling limit is that the Collect is a cosmic and ecclesial petition, not a slogan for personal risk, and turmoil is never by itself evidence of faithful direction.

\Needspace{8\baselineskip}
\noindent\textbf{Two burdens require opposite diagnoses.\enspace\properrefs{Ep., Gosp., Postcomm.}}
Romans gives creation's pain the temporal shape of labor awaiting birth; Luke gives abundance the physical shape of nets tearing and boats nearly sinking; after Communion the Church still asks for purification and protection. The mechanism is shared bodily strain under different causal and eschatological horizons. A separate reading can leave ``burden'' morally undifferentiated, whereas the conjunction asks whether a community is suffering corruption that must be resisted, fruitful increase that requires enlarged capacity, or a mixture needing purification. Its intellectual and pastoral fruit is a diagnostic discipline that refuses both to romanticize exhaustion and to discard difficult fruit too quickly. The strongest limit is that Paul's travail promises cosmic liberation while Luke narrates one apostolic sign; neither suffering nor institutional overload proves fecundity, and the Postcommunion supplies no retrospective diagnosis of the nets.

\Needspace{8\baselineskip}
\noindent\textbf{Light exposes before it steadies.\enspace\properrefs{Int., Gosp., Off.}}
The Introit first confesses the Lord as light; the catch then discloses not only fish but Peter's sinfulness and fear; at the offering the Church asks that her eyes be illuminated lest she sleep in death. The mechanism is a progression from received illumination, through disruptive self-disclosure, to a renewed petition for sight. Taken separately, confidence, confession, and vigilance can appear as unrelated moods; together they propose that grace-given light sometimes destabilizes false confidence before Christ turns fear toward vocation. The spiritual fruit is an examen that asks not only whether one has seen something true, but what that seeing reveals, what fear it produces, and whose word redirects it. The proposal's limit is that Peter's reaction is not an exposition of either psalm, and not every fear, shame, or psychological disturbance is illumination.

\Needspace{8\baselineskip}
\noindent\textbf{A catch that creates followers.\enspace\properrefs{Coll., Gosp., Sec.}}
The Gospel's command, consent, catch, confession, promise, and leaving of the boats form an agency sequence; the Collect asks that the world's course be directed, while the Secret asks God mercifully to compel rebellious wills toward himself. The proposed mechanism is divine efficacy disclosed not as inert capture but as action that elicits truthful objection, obedience, repentance, cooperation, and a new walk. Reading the net and \latin{compelle} separately can miss the question whether being powerfully drawn by grace might enlarge rather than erase creaturely agency. The theological fruit is a model of mission in which efficacy is measured by freed response, not managed compliance. Its strongest limit is decisive: Luke does not interpret the Secret or settle a metaphysics of grace, and neither net nor \latin{compelle} licenses psychological, ecclesial, or civil coercion.

\Needspace{8\baselineskip}
\noindent\textbf{A visible result still awaits a true word.\enspace\properrefs{Gosp., Off., Comm.}}
The Offertory imagines the enemy naming a fall as victory, \latin{Praevalui adversus eum}; Peter names an overwhelming success by confessing, \latin{homo peccator sum}; Christ then says \latin{Noli timere} and names a future vocation; Communion finally names the Lord as strength and refuge. The mechanism is a contest of interpretations around visible outcomes, resolved neither by the enemy's verdict nor by the recipient's self-assessment alone. Element-by-element reading misses how danger and abundance can both be misread until judgment is received from a trustworthy speaker. The fruit is a communal discipline of refusing to identify success, failure, fear, or self-condemnation with God's final word. The disconfirming condition is that these speakers do not address one event: the proposal is a liturgical comparison of speech acts, not a concealed narrative parallel or a denial of prudent empirical judgment.
